\documentclass[11pt, letterpaper]{article}

\usepackage[reals]{layout} %provide \layout macro to show current page layout

\usepackage{enumitem}

\usepackage{tabulary}

\usepackage{graphicx}
\usepackage{subfigure}
\usepackage{wrapfig}

\usepackage{amsmath}            % more evironment
\usepackage{amssymb}            % more symbol

\usepackage{ifpdf}
\ifpdf
\usepackage{epstopdf} % must be put after graphicx
\usepackage[usenames,dvipsnames]{color}
\usepackage[pdftex,bookmarks=true]{hyperref}
\hypersetup{
  pdfauthor = {Jing Liu},
  pdftitle = {Syllabus for USD courses},
  pdfsubject = {Physics},
  pdfkeywords = {nuclear,particle},
}
\pdfadjustspacing=1 %force pdfLaTeX to use the same spacing as LaTeX
\else
\usepackage[usenames,dvips]{color}
\usepackage[ps2pdf]{hyperref}
\fi

\setlength{\oddsidemargin}{0.8cm}
\setlength{\evensidemargin}{0cm}
\setlength{\textwidth}{15cm}
\setlength{\textheight}{21cm}
\setlength{\hoffset}{0cm}
\setlength{\voffset}{0cm}

% Alter some LaTeX defaults for better treatment of floats: See p.105 % of "TeX Unbound" for suggested values. See pp. 199-200 of Lamport's % "LaTeX" book for details.

% General parameters, for ALL pages:
\renewcommand{\topfraction}{0.9} % max fraction of floats at top
\renewcommand{\bottomfraction}{0.9} % max fraction of floats at bottom

% Parameters for TEXT pages (not float pages):
\setcounter{topnumber}{2}
\setcounter{bottomnumber}{2}
\setcounter{totalnumber}{4} % 2 may work better
\renewcommand{\textfraction}{0.1} % allow minimal text w. figs

% Parameters for FLOAT pages (not text pages):
\renewcommand{\floatpagefraction}{0.7}	% require fuller float pages
% N.B.: floatpagefraction MUST be less than topfraction !!

% remember to use [htp] or [htpb] for placement
%\pagestyle{headings}


\begin{document}

% -------------------------------------------------------- 
% title 
% -------------------------------------------------------- 
\begin{figure}[t!]\centering
  \includegraphics[width=0.2\textwidth]{usd.logo.eps}\hfil
  \begin{minipage}[b]{0.75\linewidth}\centering
    \rule[-4.2mm]{\textwidth}{0.3mm}\\\rule{0.98\textwidth}{0.1mm}\\
    GERDA Scientific/Technical Report Munich\\\vspace{2mm}
    GSTR-yy-Mnum \\\vspace{2mm}
    Created: February 12, 2009\\\vspace{-7.5mm}
    \rule[-5mm]{0.98\textwidth}{0.1mm}\\\rule{\textwidth}{0.3mm}\\
  \end{minipage}%
\end{figure}

\section{Page layout}
\label{sec:manual}

\layout

\section{Course Description}
\label{s:desc}

\subsection{Catalog Description}
\label{ss:cata}

This is a one-semester course in nuclear and elementary particle physics.  
Nuclear physics topics may include nuclear structure (nuclear form factors, 
multipole moments, liquid and shell models); nuclear decay; nuclear reactions; 
and other topics of interest.  Elementary particle physics topics may include 
the role of symmetry in particle physics, Quantum Electrodynamics and Quantum 
Chromodynamics; the Standard Model of Particle Physics; Strong and Weak 
interactions; Accelerator and Experimental Particle Physics; and other selected 
topics beyond the Standard Model.

\subsection{Additional Description}
\label{ss:addi}


\subsection{Description of Instructional Methods}
\label{ss:meth}


\subsection{Required Textbooks and Other Materials}
\label{ss:text}

There is no required textbook. However, the following references are 
recommended:

\begin{enumerate}[label=\arabic*]
\item Glenn F Knoll, Radiation Detection and Measurement, 3rd edition, John 
  Wiley \& Sons, Inc.(1999)
\item 
  \href{https://www.hamamatsu.com/resources/pdf/etd/PMT_handbook_v3aE.pdf}{Hamamatsu, 
  Photomultiplier Tubes – Basics and Applications, 3rd edition}
\item Helmuth Spieler, Introduction to Radiation Detectors and Electronics, 
  \href{http://www-physics.lbl.gov/~spieler/physics_198_notes/}{Lecture Notes}- 
  Physics 198, Spring Semester 1998, UC Berkeley
\end{enumerate}

\subsection{Class Attendance Policy}
\label{ss:atte}

An absence without explanation will prevent the student to get an A. Frequent 
absence will lead to a drop-off from the class.

\subsection{Course Goals}
\label{ss:goal}


\section{Evaluation Procedures}
\label{s:eval}

\subsection{Assessment}
\label{s:asse}

\subsection{Requirement on Homework}
\label{s:home}


\subsection{Performance standards/grading policy}
\label{s:grad}

There are 100 points in total, 70 for the lab activities and logs, 30 for the
final exam.

\begin{table} [htbp]
  \centering
  \label{t:grad}
  \begin{tabulary}{\linewidth}{r|l}\hline
    Grading & Percentage \\\hline
    A & 90-100\%\\
    B & 80-90\%\\
    C & 70-80\%\\
    D & 60-70\%\\
    F & $<60$\%\\
  \hline
  \end{tabulary}
\end{table}


\section{Academic Integrity}
\label{s:inte}

The College of Arts and Sciences considers plagiarism, cheating, and other
forms of academic dishonesty inimical to the objectives of higher education.
The College supports the imposition of penalties on students who engage in
academic dishonesty, as defined in the “Conduct” section of the University of
South Dakota Student Handbook.

No credit can be given for a dishonest assignment. A student found to have
engaged in any form of academic dishonesty may, at the discretion of the
instructor, be:

\begin{enumerate}[label=\textbf{\alph*.}, leftmargin=2cm]
\item Given a zero for that assignment.
\item Allowed to rewrite and resubmit the assignment for credit.
\item Assigned a reduced grade for the course.
\item Dropped from the course.
\item Failed in the course.
\end{enumerate}

\section{Disability Accommodation}
\label{s:disa}

Any student who feels s/he may need academic accommodations or access
accommodations based on the impact of a documented disability should contact
and register with Disability Services during the first week of class or as soon
as possible after the diagnosis of a disability.  Disability Services is the
official office to assist students through the process of disability
verification and coordination of appropriate and reasonable accommodations.
Students currently registered with Disability Services must obtain a new
accommodation memo each semester. 

Please note: if your home institution is not the University of South Dakota but
one of the other South Dakota Board of Regents institutions (e.g., SDSU, SDSMT,
BHSU, NSU, DSU), you should work with the disability services coordinator at
your home institution.

\vspace{0.5cm}
\begin{minipage}[b]{0.75\linewidth}
\rule{\textwidth}{0.1mm}\\
Ernetta L. Fox, Director\\
Disability Services, Room 119 Service Center\\
(605) 677-6389\\
Web Site: \href{http://www.usd.edu/ds}{www.usd.edu/ds}\\
E-mail: \href{mailto:disabilityservices@usd.edu}{disabilityservices@usd.edu}\\
\rule{\textwidth}{0.1mm}
\end{minipage}

\section{Freedom in Learning}
\label{s:free}

Under Board of Regents and University policy student academic performance may
be evaluated solely on an academic basis, not on opinions or conduct in matters
unrelated to academic standards. Students should be free to take reasoned
exception to the data or views offered in any course of study and to reserve
judgment about matters of opinion, but they are responsible for learning the
content of any course of study for which they are enrolled. Students who
believe that an academic evaluation reflects prejudiced or capricious
consideration of student opinions or conduct unrelated to academic standards
should contact the dean of the college or school that offers the class to
initiate a review of the evaluation.

\section{Diversity and Inclusive Excellence}
\label{s:dive}

The University of South Dakota strives to foster a globally inclusive learning
environment where opportunities are provided for diversity to be recognized and
respected.

\end{document}
